\section{Machine details}

Modern supercomputers are very complex and it is virtually impossible to
memorise all the hardware details of a testbed given the zoo of processor types,
generations and flavours.
While it is possible to consult your compute centre's webpages or to consult
your system administrator, it might be more reasonable to obtain that data
directly from the machine.

\begin{assessmentcrime}
 Do never grab the specifiation from the login node. Instead, ask for an
 interactive terminal on a compute node and get the data there. Many
 supercomputers' login nodes are slightly different to the compute nodes,
 i.e.~slightly different make, more memory, more cores, \ldots So you might end
 up with the wrong hardware specification if you analyse a login node.
\end{assessmentcrime}

\begin{instructions}{\linux}
 On every Linux system, information about the CPU on the system is held in a
 central file.
 You can read it out by typing \texttt{cat /proc/cpuinfo} into the terminal.
\end{instructions}


\begin{instructions}{\likwid}
 If you have Likwid installed, you can invoke \texttt{likwid-topology} to obtain
 in-depth information about your CPU.
\end{instructions}
